% GENERATED FILE. Edit canonical lessons, then run npm run generate:books.
% canonical-source-hash: fnv1a64:8d4727ab5183de24
% canonical-lessons: ML-C220-sahapravarttakan, ML-C220-udyogasthan, ML-C220-gumastan, ML-C220-enciniyar, ML-C220-vakkil

\chapter{Colleague, Official, Clerk, Engineer, Lawyer}
\label{ch:ml-colleague-official-clerk-engineer-lawyer}
\hlchaptermodality{220}

\begin{chapteropening}
\textbf{By the end of this chapter:} \emph{I can say a colleague, an official, a clerk, an engineer and a lawyer.}

Complete the last lesson of chapter 220: I can say a colleague, an official, a clerk, an engineer and a lawyer.
\end{chapteropening}

\section[\texorpdfstring{\ml{സഹപ്രവർത്തകൻ} (sahapravarttakan)}{സഹപ്രവർത്തകൻ (sahapravarttakan)}]{\ml{സഹപ്രവർത്തകൻ} (sahapravarttakan) — a colleague}

\label{lesson:ML-C220-sahapravarttakan}



\begin{quote}
Before the new one: say the Malayalam for the balance, then the Malayalam for a discount.
\end{quote}

\subsection*{You'll want to know: \ml{സഹപ്രവർത്തകൻ}}
\textbf{\ml{സഹപ്രവർത്തകൻ}} — \emph{sahapravarttakan} — \textquotedblleft{}a colleague\textquotedblright{}.

\begin{cousinweb}
\textbf{\ml{സഹപ്രവർത്തകൻ}} (\emph{sahapravarttakan}) is built from Sanskrit: \emph{saha}, \textquotedblleft{}with\textquotedblright{}, and \emph{pravartaka}, \textquotedblleft{}one who sets going, who carries on\textquotedblright{}, with the masculine ending \emph{-n}.
\end{cousinweb}

\begin{grammarlens}[title={What you've built}]
One more word for this chapter.
\end{grammarlens}

\subsection*{Your turn}
\begin{itemize}
\raggedright
  \item \emph{Say it:} \emph{sahapravarttakan}
  \item \emph{Say it:} \emph{sahapravarttakan}, once more
  \item \emph{Say it:} say \emph{sahapravarttakan}
  \item \emph{Recall:} say the Malayalam for the balance, then the Malayalam for a discount, then say \emph{sahapravarttakan} again
\end{itemize}

\begin{tcolorbox}[breakable,colback=teal!4,colframe=teal!35!black,title={Before you move on}]
What does \ml{സഹപ്രവർത്തകൻ} mean? (\textquotedblleft{}A colleague\textquotedblright{}.) Say it once more.
\end{tcolorbox}
\section[\texorpdfstring{\ml{ഉദ്യോഗസ്ഥൻ} (udyōgasthan)}{ഉദ്യോഗസ്ഥൻ (udyōgasthan)}]{\ml{ഉദ്യോഗസ്ഥൻ} (udyōgasthan) — an official, an officer}

\label{lesson:ML-C220-udyogasthan}



\begin{quote}
Before the new one: say the Malayalam for a discount, then the Malayalam for a colleague.
\end{quote}

\subsection*{You'll want to know: \ml{ഉദ്യോഗസ്ഥൻ}}
\textbf{\ml{ഉദ്യോഗസ്ഥൻ}} — \emph{udyōgasthan} — \textquotedblleft{}an official, an officer\textquotedblright{}.

\begin{cousinweb}
\textbf{\ml{ഉദ്യോഗസ്ഥൻ}} (\emph{udyōgasthan}) is built from Sanskrit \emph{udyōga}, \textquotedblleft{}exertion; employment\textquotedblright{} (CDIAL 2055), and \emph{-stha}, \textquotedblleft{}standing in\textquotedblright{}, with the masculine ending \emph{-n}.
\end{cousinweb}

\begin{grammarlens}[title={What you've built}]
One more word for this chapter.
\end{grammarlens}

\subsection*{Your turn}
\begin{itemize}
\raggedright
  \item \emph{Say it:} \emph{udyōgasthan}
  \item \emph{Say it:} \emph{udyōgasthan}, once more
  \item \emph{Say it:} say \emph{udyōgasthan}
  \item \emph{Recall:} say the Malayalam for a discount, then the Malayalam for a colleague, then say \emph{udyōgasthan} again
\end{itemize}

\begin{tcolorbox}[breakable,colback=teal!4,colframe=teal!35!black,title={Before you move on}]
What does \ml{ഉദ്യോഗസ്ഥൻ} mean? (\textquotedblleft{}An official, an officer\textquotedblright{}.) Say it once more.
\end{tcolorbox}
\section[\texorpdfstring{\ml{ഗുമസ്തൻ} (gumastan)}{ഗുമസ്തൻ (gumastan)}]{\ml{ഗുമസ്തൻ} (gumastan) — a clerk}

\label{lesson:ML-C220-gumastan}



\begin{quote}
Before the new one: say the Malayalam for a colleague, then the Malayalam for an official.
\end{quote}

\subsection*{You'll want to know: \ml{ഗുമസ്തൻ}}
\textbf{\ml{ഗുമസ്തൻ}} — \emph{gumastan} — \textquotedblleft{}a clerk\textquotedblright{}.

\begin{cousinweb}
\textbf{\ml{ഗുമസ്തൻ}} (\emph{gumastan}) $\leftarrow$ Persian \emph{gumāshta}, \textquotedblleft{}an agent, one appointed\textquotedblright{}, which came through Hindustani with the clerks of the courts.
\end{cousinweb}

\begin{grammarlens}[title={What you've built}]
One more word for this chapter.
\end{grammarlens}

\subsection*{Your turn}
\begin{itemize}
\raggedright
  \item \emph{Say it:} \emph{gumastan}
  \item \emph{Say it:} \emph{gumastan}, once more
  \item \emph{Say it:} say \emph{gumastan}
  \item \emph{Recall:} say the Malayalam for a colleague, then the Malayalam for an official, then say \emph{gumastan} again
\end{itemize}

\begin{tcolorbox}[breakable,colback=teal!4,colframe=teal!35!black,title={Before you move on}]
What does \ml{ഗുമസ്തൻ} mean? (\textquotedblleft{}A clerk\textquotedblright{}.) Say it once more.
\end{tcolorbox}
\section[\texorpdfstring{\ml{എഞ്ചിനീയർ} (eñcinīyar)}{എഞ്ചിനീയർ (eñcinīyar)}]{\ml{എഞ്ചിനീയർ} (eñcinīyar) — an engineer}

\label{lesson:ML-C220-enciniyar}



\begin{quote}
Before the new one: say the Malayalam for an official, then the Malayalam for a clerk.
\end{quote}

\subsection*{You'll want to know: \ml{എഞ്ചിനീയർ}}
\textbf{\ml{എഞ്ചിനീയർ}} — \emph{eñcinīyar} — \textquotedblleft{}an engineer\textquotedblright{}.

\begin{cousinweb}
\textbf{\ml{എഞ്ചിനീയർ}} (\emph{eñcinīyar}) $\leftarrow$ English \emph{engineer}.
\end{cousinweb}

\begin{grammarlens}[title={What you've built}]
One more word for this chapter.
\end{grammarlens}

\subsection*{Your turn}
\begin{itemize}
\raggedright
  \item \emph{Say it:} \emph{eñcinīyar}
  \item \emph{Say it:} \emph{eñcinīyar}, once more
  \item \emph{Say it:} say \emph{eñcinīyar}
  \item \emph{Recall:} say the Malayalam for an official, then the Malayalam for a clerk, then say \emph{eñcinīyar} again
\end{itemize}

\begin{tcolorbox}[breakable,colback=teal!4,colframe=teal!35!black,title={Before you move on}]
What does \ml{എഞ്ചിനീയർ} mean? (\textquotedblleft{}An engineer\textquotedblright{}.) Say it once more.
\end{tcolorbox}
\section[\texorpdfstring{\ml{വക്കീൽ} (vakkīl)}{വക്കീൽ (vakkīl)}]{\ml{വക്കീൽ} (vakkīl) — a lawyer}

\label{lesson:ML-C220-vakkil}



\begin{quote}
Before the new one: say the Malayalam for a clerk, then the Malayalam for an engineer.
\end{quote}

\subsection*{You'll want to know: \ml{വക്കീൽ}}
\textbf{\ml{വക്കീൽ}} — \emph{vakkīl} — \textquotedblleft{}a lawyer\textquotedblright{}.

\begin{cousinweb}
\textbf{\ml{വക്കീൽ}} (\emph{vakkīl}) $\leftarrow$ Arabic \emph{wakīl}, \textquotedblleft{}agent, representative\textquotedblright{}, which came through Persian and Hindustani.
\end{cousinweb}

\begin{grammarlens}[title={What you've built}]
One more word for this chapter.
\end{grammarlens}

\subsection*{Your turn}
\begin{itemize}
\raggedright
  \item \emph{Say it:} \emph{vakkīl}
  \item \emph{Say it:} \emph{vakkīl}, once more
  \item \emph{Say it:} say \emph{vakkīl}
  \item \emph{Recall:} say the Malayalam for a clerk, then the Malayalam for an engineer, then say \emph{vakkīl} again
\end{itemize}

\begin{tcolorbox}[breakable,colback=teal!4,colframe=teal!35!black,title={Before you move on}]
What does \ml{വക്കീൽ} mean? (\textquotedblleft{}A lawyer\textquotedblright{}.) Say it once more.
\end{tcolorbox}
